\begin{table}[tb]
  \centering
  \caption{Reduced critical Reynolds number against the full-order lift departure at selected amplitudes.}
  \label{tab:critcurve}
  \begin{tabular}{lrrrr}
    \toprule
    $\arot$ & $\Rec(\mu)$ & bracket & $\Rec(\CL^{\mathrm{FOM}})$ & $|\Delta|$ \\
    \midrule
    -1.00 & 40.44 & 1 & 41.00 & 0.56 \\
    -0.60 & 48.40 & 1 & 49.00 & 0.60 \\
    -0.40 & 53.75 & 1 & 89.00$^{\dagger}$ & 35.25 \\
    -0.20 & 60.27 & 1 & 61.00 & 0.73 \\
    +0.20 & 77.97 & 2 & 78.00 & 0.03 \\
    +0.60 & 105.41 & 1 & 103.00 & 2.41 \\
    \bottomrule
  \end{tabular}
  \\[2pt] \footnotesize $\Rec(\mu)$ is the zero crossing of the reduced stability indicator, bracketed by one continuation step. The reference column is full-order \emph{lift departure}, not bisection: the first $\Rey$ at which an asymmetric leg's $|\CL^{\mathrm{FOM}}|$ clears a fixed fraction of its own scale. Rows marked $\dagger$ are amplitudes at which the proxy does not locate the pitchfork because the asymmetric leg was not continued down to it: its stored full-order points begin well above $\Rec$, so the first $\Rey$ at which lift departs zero is pinned to the start of a truncated leg rather than to the bifurcation. They are shown for completeness and excluded from the quoted agreement. $\Rec(\mu)$ is taken as the \emph{first} real zero crossing. The reduced indicator is cluster-local and discontinuous where the cluster assignment changes along a leg, which introduces spurious later crossings; on this $K=4$ configuration every such crossing falls above the genuine pitchfork, so the first-crossing rule is unaffected. That ordering is a property of the cluster partition, not a guarantee: it is validated here, not robust in general. A full-order eigensweep (and a bisection driver) would replace this column outright; see \cref{sec:conclusions}.
\end{table}
